
%%=============================================================================
\chapter{Anhang}
\label{chap:anhang}
%%=============================================================================


%%=============================================================================
\section{Schaltplan}
\label{sec:APP.Sch}

\begin{table}[h]
    \centering
        \small
        \begin{tabular} {rll}  \hline
    Seite& Titel & Beschreibung \\ \hline
      1  & Forth-Board &   CPU, Port- und Signalbelegung \\
      2  & Adressdekod. D/A-Wandler, LED's & Daten- u. Adressbus, Analogausgänge \\
      3  & Analogsignalverarbeitung    &   Analogeingänge mit Schutzbeschaltung \\
      4  & Einspeisung, Reset-Signal   &       \\
      5  & Stromerfassung  &   Anschl. LEM-Wandler, INA114 \\
      6  & Toleranzbandregelung Phasenströme   &   Komparatoren \\
      7  & I1\_aus, Überspannungserkennung  & \\
      8  & Stromregellogik, Strangfreigaben    & GALs\\
      9  & Treiberstufen   &       \\
     10  &      &       \\
     11  & Messung Phasen- u. ZWK-Spannung &   U-Messung, Differenzverstärker   \\
     12  &     &       \\
     13  & Drehmomentberechnung    &   DATC \\
     14  & Momentberechnung Sample\&Hold    &   T\_Sample \\
     15  &     &       \\
     16  & Lagegebersignalverarbeitung &    \\
     17  & Temp.-Überwachung Machine + UMR &    \\
     18  & Schnittstellen: RS232, CAN &     \\  \hline
        \end{tabular}
        \caption{Übersicht über den Schaltplan \ttt{OKOFEH\_CONTROL}}
        \label{tab:Inhalt-OKOFEH-Control}
\end{table}

\newpage
\ifpdf
	\begin{landscape}		%% Landscape/Querformat im PDF-Dokument
\fi
	\centering
	\includegraphics*[height=.85\textheight]{cntr-s01}
  \label{sch:cs1}
\ifpdf
	\end{landscape}
\fi
